\chapter{Modelo Entidade-Relacionamento Conceitual (MER)}
\label{cap_modelo_conceitual}

O Modelo Entidade-Relacionamento Conceitual (MER) constitui o nível mais elevado de abstração na engenharia de dados, representando os conceitos do domínio de aplicação, seus atributos intrínsecos e as associações estruturais de forma independente de qualquer consideração de implementação física de software ou Sistema de Gerenciamento de Banco de Dados \cite{chen1976, elmasri2016}.

No SIGEA-GTT, a modelagem conceitual sintetiza a interação entre os processos pedagógicos universitários e o método internacional de rastreamento de eventos adversos hospitalares.

% ===========================================================================
\section{Identificação e Classificação das Entidades do Domínio}
\label{sec_entidades_conceituais}
% ===========================================================================

As entidades identificadas no domínio são categorizadas formalmente de acordo com sua autonomia de existência e dependência estrutural:

\subsection{Entidades Fortes (Regulares ou Independentes)}
Possuem existência autônoma no minimundo, com chaves identificadoras intrínsecas e independentes da existência de outras entidades:
\begin{enumerate}
    \item \textbf{USUÁRIO}: Representa os atores autenticáveis no ecossistema institucional. Armazena o identificador único universal, o nome completo, o correio eletrônico institucional sob o domínio exclusivo \texttt{@academico.ufs.br}, o resumo criptográfico da senha (\textit{hash}), o papel de acesso (\texttt{ADMIN}, \texttt{PROFESSOR} ou \texttt{STUDENT}), o número de matrícula acadêmica (restrito exclusivamente a discentes), o estado de ativação de conta e os sinalizadores de obrigatoriedade de redefinição de credencial provisória;
    \item \textbf{MÓDULO\_GTT}: Representa cada uma das seis áreas clínicas de investigação sistemática preconizadas pelo \textit{Institute for Healthcare Improvement} (IHI). Armazena o identificador único, o código mnemônico internacional (\texttt{C}, \texttt{S}, \texttt{M}, \texttt{I}, \texttt{P}, \texttt{E}), a denominação do módulo e o escopo assistencial coberto;
    \item \textbf{GRAVIDADE\_DANO}: Representa a taxonomia de severidade de dano ao paciente estabelecida pelo conselho norte-americano NCC MERP. Armazena a letra da categoria (de \texttt{A} a \texttt{I}), a denominação clínica, a descrição do grau de intervenção exigido e o atributo booleano de dano real, que segrega os incidentes que não causaram dano (\texttt{A}--\texttt{D}) dos eventos adversos com dano manifesto (\texttt{E}--\texttt{I});
    \item \textbf{TURMA\_ACADÊMICA}: Representa a instância semestral de uma disciplina letiva ofertada no curso de Enfermagem da UFS. Armazena o nome da matéria curricular, o código da turma, o período letivo formatado, a situação de encerramento da turma e a data de criação.
\end{enumerate}

\subsection{Entidades Fracas e Dependentes por Existência}
Dependem existencialmente de entidades fortes para ter significado no modelo:
\begin{enumerate}
    \item \textbf{GATILHO\_CLÍNICO}: Representa cada um dos 53 rastreadores objetivos de eventos adversos (ex.: \textit{"Uso de Naloxone", "Queda do Paciente", "Reoperação não planejada"}). Depende de \textbf{MÓDULO\_GTT}, possuindo código alfanumérico identificador, nome clínico, diretrizes de busca em prontuário e indicador ativo;
    \item \textbf{ATIVIDADE\_AUDITORIA}: Representa a unidade didática proposta por um docente para uma turma específica. Depende de \textbf{TURMA\_ACADÊMICA}, armazenando título, orientações, prazo de entrega e o prontuário simulado completo em alta fidelidade;
    \item \textbf{SUBMISSÃO\_AVALIAÇÃO}: Representa a resolução entregue por um discente. Depende conjuntamente de \textbf{ATIVIDADE\_AUDITORIA} e do \textbf{USUÁRIO} discente, armazenando os gatilhos detectados com justificativa clínica, ferramentas de qualidade, data de envio, nota e parecer docente.
\end{enumerate}

\subsection{Entidade Associativa}
\begin{enumerate}
    \item \textbf{MATRÍCULA\_DISCENTE}: Modela o relacionamento muitos-para-muitos (N:M) entre discentes e turmas acadêmicas. Concretiza o ato acadêmico de inscrição de um \textbf{USUÁRIO} discente em uma \textbf{TURMA\_ACADÊMICA}, possuindo como atributo próprio a data/hora de efetivação da matrícula.
\end{enumerate}

% ===========================================================================
\section{Definição dos Relacionamentos e Cardinalidades}
\label{sec_relacionamentos_mer}
% ===========================================================================

Os vínculos conceituais entre as entidades, suas cardinalidades mínimas e máximas $(min, max)$, e seus significados semânticos são especificados a seguir:

\begin{itemize}
    \item \textbf{Ministra}: Associação 1:N entre \textbf{USUÁRIO} (Docente) e \textbf{TURMA\_ACADÊMICA}.
    \begin{itemize}
        \item Cardinalidade de \textbf{USUÁRIO}: $(0, N)$ --- Um usuário docente pode ministrar nenhuma ou várias turmas acadêmicas;
        \item Cardinalidade de \textbf{TURMA\_ACADÊMICA}: $(1, 1)$ --- Cada turma acadêmica possui obrigatoriamente um e apenas um professor titular responsável.
    \end{itemize}

    \item \textbf{Inscrição}: Associação N:M entre \textbf{USUÁRIO} (Discente) e \textbf{TURMA\_ACADÊMICA}, mediada por \textbf{MATRÍCULA\_DISCENTE}.
    \begin{itemize}
        \item Cardinalidade de \textbf{USUÁRIO}: $(0, N)$ --- Um discente pode estar matriculado em zero ou várias turmas no semestre;
        \item Cardinalidade de \textbf{TURMA\_ACADÊMICA}: $(0, N)$ --- Uma turma pode conter de zero a múltiplos alunos inscritos.
    \end{itemize}

    \item \textbf{Compõe}: Associação 1:N entre \textbf{MÓDULO\_GTT} e \textbf{GATILHO\_CLÍNICO}.
    \begin{itemize}
        \item Cardinalidade de \textbf{MÓDULO\_GTT}: $(1, N)$ --- Um módulo do IHI é composto por um ou múltiplos gatilhos clínicos padronizados;
        \item Cardinalidade de \textbf{GATILHO\_CLÍNICO}: $(1, 1)$ --- Cada gatilho pertence estritamente a um único módulo de auditoria.
    \end{itemize}

    \item \textbf{Disponibiliza}: Relação 1:N entre turma e atividade acadêmica;
    \begin{itemize}
        \item Cardinalidade de \textbf{TURMA\_ACADÊMICA}: $(0, N)$ --- Uma turma pode possuir zero ou diversas atividades;
        \item Cardinalidade de \textbf{ATIVIDADE\_AUDITORIA}: $(1, 1)$ --- Cada atividade vincula-se a exatamente uma turma.
    \end{itemize}

    \item \textbf{Recebe Submissão}: Relação 1:N entre atividade e submissões discentes;
    \begin{itemize}
        \item Cardinalidade de \textbf{ATIVIDADE\_AUDITORIA}: $(0, N)$ --- Uma atividade recebe de zero a múltiplas submissões;
        \item Cardinalidade de \textbf{SUBMISSÃO\_AVALIAÇÃO}: $(1, 1)$ --- Cada submissão pertence a uma única atividade.
    \end{itemize}

    \item \textbf{Submete}: Associação 1:N entre usuário discente e submissões entregues;
    \begin{itemize}
        \item Cardinalidade de \textbf{USUÁRIO}: $(0, N)$ --- Um discente realiza zero ou múltiplas submissões;
        \item Cardinalidade de \textbf{SUBMISSÃO\_AVALIAÇÃO}: $(1, 1)$ --- Cada submissão é de autoria de um único discente.
    \end{itemize}
\end{itemize}

% ===========================================================================
\section{Diagrama Entidade-Relacionamento Conceitual (Notação de Chen)}
\label{sec_diagrama_mer_chen}
% ===========================================================================

A Figura \ref{fig_mer_conceitual} ilustra o Modelo Entidade-Relacionamento Conceitual integral do SIGEA-GTT fundamentado na notação formal de Peter Chen \cite{chen1976}, utilizando retângulos para entidades fortes, retângulos duplos para entidades fracas, losangos para relacionamentos normais, losangos duplos para relacionamentos identificadores e elipses para atributos chave e descritivos.

\begin{figure}[!ht]
\centering
\caption{Modelo Entidade-Relacionamento Conceitual (Notação Canônica de Peter Chen)}
\label{fig_mer_conceitual}
\resizebox{0.98\textwidth}{!}{%
\begin{tikzpicture}[
    node distance=1.8cm and 2.2cm,
    every text node part/.style={align=center},
    % Estilos das Entidades
    entity/.style={draw=blue!80!black, fill=blue!5, thick, rectangle, minimum width=3.2cm, minimum height=1.1cm, font=\small\bfseries},
    weak entity/.style={draw=blue!80!black, fill=blue!5, thick, double, rectangle, minimum width=3.2cm, minimum height=1.1cm, font=\small\bfseries},
    assoc entity/.style={draw=teal!80!black, fill=teal!5, thick, rectangle, minimum width=3.4cm, minimum height=1.1cm, font=\small\bfseries},
    % Estilos dos Relacionamentos
    relationship/.style={draw=gray!80!black, fill=gray!10, thick, diamond, aspect=1.8, minimum width=2.4cm, minimum height=1.1cm, font=\scriptsize\bfseries},
    ident relationship/.style={draw=gray!80!black, fill=gray!10, thick, double, diamond, aspect=1.8, minimum width=2.4cm, minimum height=1.1cm, font=\scriptsize\bfseries},
    % Estilos dos Atributos
    attribute/.style={draw=black!70, fill=white, thin, ellipse, minimum width=1.8cm, minimum height=0.7cm, font=\tiny},
    key attribute/.style={draw=black!70, fill=gray!20, thick, ellipse, minimum width=1.8cm, minimum height=0.7cm, font=\tiny\bfseries},
    edge/.style={draw=black!80, semithick}
]

    % ==========================================
    % LINHA 1: USUÁRIO E TURMA ACADÊMICA
    % ==========================================
    \node[entity] (usuario) at (0, 8) {USUÁRIO};
    \node[relationship] (ministra) at (5, 8) {MINISTRA};
    \node[entity] (turma) at (10, 8) {TURMA\\ACADÊMICA};

    % Atributos do USUÁRIO
    \node[key attribute] (u_id) at (-2.4, 9.5) {\underline{id (UUID)}};
    \node[attribute] (u_nome) at (-2.4, 8.5) {nome\_completo};
    \node[attribute] (u_email) at (-2.4, 7.5) {email};
    \node[attribute] (u_role) at (-1.2, 6.5) {role};
    \node[attribute] (u_matr) at (1.2, 6.5) {matricula};

    \draw[edge] (usuario) -- (u_id);
    \draw[edge] (usuario) -- (u_nome);
    \draw[edge] (usuario) -- (u_email);
    \draw[edge] (usuario) -- (u_role);
    \draw[edge] (usuario) -- (u_matr);

    % Atributos da TURMA
    \node[key attribute] (t_id) at (12.5, 9.5) {\underline{id (UUID)}};
    \node[attribute] (t_mat) at (12.5, 8.5) {nome\_materia};
    \node[attribute] (t_cod) at (12.5, 7.5) {codigo\_turma};
    \node[attribute] (t_per) at (11.5, 6.5) {periodo\_letivo};

    \draw[edge] (turma) -- (t_id);
    \draw[edge] (turma) -- (t_mat);
    \draw[edge] (turma) -- (t_cod);
    \draw[edge] (turma) -- (t_per);

    % Conexões Ministro
    \draw[edge] (usuario) -- node[above, font=\scriptsize] {(0, N)} (ministra);
    \draw[edge] (ministra) -- node[above, font=\scriptsize] {(1, 1)} (turma);

    % ==========================================
    % ASSOCIAÇÃO MATRÍCULA (N:M)
    % ==========================================
    \node[relationship] (matricula_rel) at (5, 6) {CURSA};
    \node[assoc entity] (matricula) at (5, 4.4) {MATRÍCULA\\DISCENTE};
    \node[attribute] (m_data) at (5, 3.2) {data\_matricula};

    \draw[edge] (usuario.south) |- node[near end, left, font=\scriptsize] {(0, N)} (matricula_rel.west);
    \draw[edge] (turma.south) |- node[near end, right, font=\scriptsize] {(0, N)} (matricula_rel.east);
    \draw[edge, dashed] (matricula_rel) -- (matricula);
    \draw[edge] (matricula) -- (m_data);

    % ==========================================
    % LINHA 2: ATIVIDADE E SUBMISSÃO
    % ==========================================
    \node[ident relationship] (propoe) at (10, 3) {DISPONIBILIZA};
    \node[weak entity] (atividade) at (10, 0.8) {ATIVIDADE\\AUDITORIA};

    \draw[edge] (turma) -- node[right, font=\scriptsize] {(1, 1)} (propoe);
    \draw[edge] (propoe) -- node[right, font=\scriptsize] {(0, N)} (atividade);

    % Atributos da ATIVIDADE
    \node[key attribute] (a_id) at (12.8, 1.8) {\underline{id (UUID)}};
    \node[attribute] (a_tit) at (12.8, 0.8) {titulo};
    \node[attribute] (a_pront) at (12.8, -0.2) {prontuario (JSONB)};
    \node[attribute] (a_prazo) at (10, -0.7) {data\_limite};

    \draw[edge] (atividade) -- (a_id);
    \draw[edge] (atividade) -- (a_tit);
    \draw[edge] (atividade) -- (a_pront);
    \draw[edge] (atividade) -- (a_prazo);

    % SUBMISSÃO
    \node[ident relationship] (recebe) at (5, 0.8) {RECEBE};
    \node[weak entity] (submissao) at (0, 0.8) {SUBMISSÃO\\AVALIAÇÃO};

    \draw[edge] (atividade) -- node[above, font=\scriptsize] {(0, N)} (recebe);
    \draw[edge] (recebe) -- node[above, font=\scriptsize] {(1, 1)} (submissao);

    % Vínculo Discente -> Submissão
    \node[relationship] (submete_rel) at (0, 3.8) {SUBMETE};
    \draw[edge] (usuario) -- node[left, font=\scriptsize] {(0, N)} (submete_rel);
    \draw[edge] (submete_rel) -- node[left, font=\scriptsize] {(1, 1)} (submissao);

    % Atributos da SUBMISSÃO
    \node[key attribute] (s_id) at (-2.5, 1.8) {\underline{id (UUID)}};
    \node[attribute] (s_gat) at (-2.5, 0.8) {gatilhos (JSONB)};
    \node[attribute] (s_ferr) at (-2.5, -0.2) {ferramentas (JSONB)};
    \node[attribute] (s_nota) at (0, -0.7) {nota / parecer};

    \draw[edge] (submissao) -- (s_id);
    \draw[edge] (submissao) -- (s_gat);
    \draw[edge] (submissao) -- (s_ferr);
    \draw[edge] (submissao) -- (s_nota);

    % ==========================================
    % TAXONOMIA CLÍNICA: MÓDULO E GATILHOS
    % ==========================================
    \node[entity] (modulo) at (-2, -3.5) {MÓDULO\\GTT};
    \node[ident relationship] (compoe) at (3.5, -3.5) {COMPÕE};
    \node[weak entity] (gatilho) at (8.5, -3.5) {GATILHO\\CLÍNICO};

    \draw[edge] (modulo) -- node[above, font=\scriptsize] {(1, 1)} (compoe);
    \draw[edge] (compoe) -- node[above, font=\scriptsize] {(1, N)} (gatilho);

    % Atributos Módulo
    \node[key attribute] (m_id_gtt) at (-4, -2.6) {\underline{id (UUID)}};
    \node[attribute] (m_cod) at (-4, -3.5) {codigo};
    \node[attribute] (m_nome) at (-4, -4.4) {nome};

    \draw[edge] (modulo) -- (m_id_gtt);
    \draw[edge] (modulo) -- (m_cod);
    \draw[edge] (modulo) -- (m_nome);

    % Atributos Gatilho
    \node[key attribute] (g_id) at (10.6, -2.6) {\underline{id (UUID)}};
    \node[attribute] (g_cod) at (10.6, -3.5) {codigo};
    \node[attribute] (g_nome) at (10.6, -4.4) {nome / descricao};

    \draw[edge] (gatilho) -- (g_id);
    \draw[edge] (gatilho) -- (g_cod);
    \draw[edge] (gatilho) -- (g_nome);

    % GRAVIDADE DE DANO (Catálogo Canônico)
    \node[entity] (gravidade) at (3.5, -5.5) {GRAVIDADE\\DANO};
    \node[key attribute] (gr_id) at (1.2, -6.5) {\underline{id (UUID)}};
    \node[attribute] (gr_let) at (3.5, -6.8) {categoria (A--I)};
    \node[attribute] (gr_dano) at (5.8, -6.5) {is\_harm};

    \draw[edge] (gravidade) -- (gr_id);
    \draw[edge] (gravidade) -- (gr_let);
    \draw[edge] (gravidade) -- (gr_dano);

\end{tikzpicture}
}
\fonte{Elaborado pelo autor (2026).}
\end{figure}
